%% raw_s3a.tex -- Slide Set 3 ("Logistic models", Week 3), PDF pages 1-42 (the tested part).
%% Raw transcription, trimmed: only material not already in Part I is kept.
\section{Slide Set 3 --- Logistic regression \& model evaluation (pp.~1--42, the tested part)}
\ixh{logistic regression}\ixh{model evaluation}
{\color{figgray}\itshape [Slides/parts already covered in Part I are omitted: S3.1--4, S3.6--7, S3.9--25, S3.28--29,
S3.32--34, S3.37--38, S3.41--42; on the kept slides only the details missing from Part I remain.]}

\subsection{Overview: classification}
\ixh{classification}

%% ---------------------------------------------------------------- S3.5
\slide{3}{5}{Why not Linear Regression?}
\ixh{linear regression}\ixh{linear probability model}
\fig{\emph{Probability of Default} vs.\ \emph{Balance} (0--2500). \textbf{Left, ``Linear Regression''}:
the fitted blue straight line starts below zero (about $-0.1$ at balance 0), crosses 0 near balance 500--600 and
only reaches about 0.28 at balance 2700. \textbf{Right, ``Logistic Regression''}: the blue fitted
curve is S-shaped (\ix{sigmoid}): essentially 0 up to balance $\approx1000$, rising steeply between $\approx1500$
and $\approx2200$ (0.5 at about 1950--2000), and levelling off near 0.98 at 2650.}

%% ---------------------------------------------------------------- S3.8
\slide{3}{8}{A data example}
\ixh{Default data set}
\fig{three panels (the ISLR \texttt{Default} data). \textbf{Left}: \ix{scatterplot} of \emph{Income} (y, 0 to
$\approx70000$) vs.\ \emph{Balance} (x, 0 to $\approx2700$); the defaulters (orange ``+'') sit almost entirely at
balance $\gtrsim1200$, most of them 1500--2500, at every income level. \textbf{Middle}: \ix{boxplot}s of
\emph{Balance} by \emph{Default}: No has median $\approx800$, box $\approx470$--$1130$, whiskers 0 to $\approx2120$
and high outliers up to $\approx2390$; Yes has median $\approx1790$, box $\approx1510$--$1990$, whiskers
$\approx960$--$2650$ and low outliers at $\approx650$--$780$. \textbf{Right}: boxplots of \emph{Income} by
\emph{Default}: No median $\approx34500$ (box $\approx21500$--$44000$, whiskers $\approx800$--$73500$), Yes median
$\approx31500$ (box $\approx19000$--$43000$, whiskers $\approx10000$--$66500$).}

\subsection{Estimating the coefficients}
\ixh{maximum likelihood}

%% ---------------------------------------------------------------- S3.26
\slide{3}{26}{Estimating coefficients}
\textbf{Maximum Likelihood Estimation (MLE)}\ixh{maximum likelihood}\ixh{MLE}
\begin{routput}
Call:
glm(formula = default ~ student + balance + income,  (*@\ixch{glm()}\ixch{Default}@*)
    family = binomial, data = Default)  (*@\ixch{family = binomial}@*)

Deviance Residuals:
    Min       1Q   Median       3Q      Max
-2.4691  -0.1418  -0.0557  -0.0203   3.7383

Coefficients:
...
---
Signif. codes:
0 '***' 0.001 '**' 0.01 '*' 0.05 '.' 0.1 ' ' 1

(Dispersion parameter for binomial family taken to be 1)
...
Number of Fisher Scoring iterations: 8
\end{routput}
{\color{figgray}\itshape [note: the elided rows (coefficient table, null/residual deviance, AIC) are in Part I.]}

\subsection{Confounding}
\ixh{confounding}

%% ---------------------------------------------------------------- S3.27
\slide{3}{27}{What is wrong?}
A simple logistic regression of \textbf{default} on \textbf{student}\ixh{simple logistic regression}

\begin{tabular}{@{}l r r r r@{}}
\toprule
 & Coefficient & Std.\ error & Z-statistic & P-value\\
\midrule
\texttt{\textcolor{rawc}{Intercept}} & $-3.5041$ & 0.0707 & $-49.55$ & $<0.0001$\\
\texttt{\textcolor{rawc}{student[Yes]}} & \fbox{0.4049} & 0.1150 & 3.52 & 0.0004\\
\bottomrule
\end{tabular}
\ixh{coefficient}\ixh{standard error}\ixh{z-statistic}\ixh{p-value}

A multiple logistic regression of \textbf{default} on \textbf{balance}, \textbf{income}, and \textbf{student}\ixh{multiple logistic regression}

\begin{tabular}{@{}l r r r r@{}}
\toprule
 & Coefficient & Std.\ error & Z-statistic & P-value\\
\midrule
\texttt{\textcolor{rawc}{Intercept}} & $-10.8690$ & 0.4923 & $-22.08$ & $<0.0001$\\
\texttt{\textcolor{rawc}{balance}} & 0.0057 & 0.0002 & 24.74 & $<0.0001$\\
\texttt{\textcolor{rawc}{income}} & 0.0030 & 0.0082 & 0.37 & 0.7115\\
\texttt{\textcolor{rawc}{student[Yes]}} & \fbox{$-0.6468$} & 0.2362 & $-2.74$ & 0.0062\\
\bottomrule
\end{tabular}

%% ---------------------------------------------------------------- S3.30
\slide{3}{30}{[Simpson's paradox example]}
\ixh{Simpson's paradox}
\fig{Blue group: points $(1,6),(2,7),(3,8),(4,9)$ with a solid blue line of slope $+1$ through them. Red group: points
$(8,1),(9,2),(10,3),(11,4)$ with a solid red line of slope $+1$. A black dashed line fitted to all eight points
together slopes \emph{downward}, from $y\approx9.2$ at the left edge ($x\approx-1.5$) to $y\approx0.8$ at the right
edge ($x\approx13.5$), i.e.\ slope $\approx-0.56$.}

%% ---------------------------------------------------------------- S3.31
\slide{3}{31}{[Simpson's paradox: scatter with clusters]}
\ixh{Simpson's paradox}\ixh{correlation}
\textsf{Korrelation:}
\fig{a ggplot-style \ix{scatterplot} (title ``Korrelation:'' with no value shown --- a frame of an
animation) of about 500 black points, $Y$ (vertical, ticks 0, 5, 10, 15) vs.\ $X$ (horizontal, ticks 0, 5, 10).
The points form about five overlapping clusters arranged along a diagonal from the upper left ($X\approx1$--$5$,
$Y\approx7$--$13$) to the lower right ($X\approx9$--$13$, $Y\approx1$--$6$), so the overall correlation of $X$ and
$Y$ is clearly negative, while inside each cluster the points tilt upward (positive within-cluster relationship).}

\subsection{Classification performance: confusion matrix}
\ixh{model evaluation}\ixh{confusion matrix}

%% ---------------------------------------------------------------- S3.35
\slide{3}{35}{Confusion matrices (multiclass)}
\ixh{confusion matrix}\ixh{multiclass classification}

\begin{tabular}{@{}l l c c c@{}}
 & & \multicolumn{3}{c}{\textcolor{OliveGreen}{\textbf{True Class}}}\\
 & & Apple & Orange & Mango\\
\cmidrule(l){3-5}
\multirow{3}{*}{\rotatebox{90}{\textcolor{Salmon}{\textbf{Predicted}}}} & Apple & \cellcolor{green!20}7 & \cellcolor{red!25}8 & \cellcolor{red!25}9\\
 & Orange & \cellcolor{red!25}1 & \cellcolor{green!20}2 & \cellcolor{red!25}3\\
 & Mango & \cellcolor{red!25}3 & \cellcolor{red!25}2 & \cellcolor{green!20}1\\
\end{tabular}

%% ---------------------------------------------------------------- S3.36
\slide{3}{36}{Many definitions and synonyms}
\ixh{confusion matrix}
A large table (Wikipedia ``Receiver operating characteristic''): \url{https://en.wikipedia.org/wiki/Receiver_operating_characteristic}

\textbf{Row ratios} (right of the block):
\begin{itemize}
\item \textbf{False discovery rate (FDR)} $=\dfrac{\Sigma\,\text{False positive}}{\Sigma\,\text{Predicted condition positive}}$\ixh{false discovery rate (FDR)}
\item \textbf{False omission rate (FOR)} $=\dfrac{\Sigma\,\text{False negative}}{\Sigma\,\text{Predicted condition negative}}$\ixh{false omission rate (FOR)}
\end{itemize}
\textbf{Ratios of ratios} (bottom right):
\begin{itemize}
\item \textbf{Positive likelihood ratio (LR+)} $=\dfrac{\text{TPR}}{\text{FPR}}$; \quad
\textbf{Negative likelihood ratio (LR$-$)} $=\dfrac{\text{FNR}}{\text{TNR}}$\ixh{likelihood ratio}
\item \textbf{Diagnostic odds ratio (DOR)} $=\dfrac{\text{LR+}}{\text{LR}-}$\ixh{diagnostic odds ratio}
\end{itemize}

\subsection{ROC curve and AUC}
\ixh{ROC curve}\ixh{AUC}

%% ---------------------------------------------------------------- S3.39
\slide{3}{39}{ROC}
\ixh{ROC curve}
\fig{plot titled ``ROC Curve'': the blue step-like ROC curve rises almost vertically
from $(0,0)$ to TPR $\approx0.5$ at FPR $\approx0.01$, reaches $\approx0.8$ at FPR $\approx0.07$, $\approx0.9$ at
FPR $\approx0.15$, $\approx0.95$ at FPR $\approx0.25$, $\approx0.99$ by FPR $\approx0.4$ and 1.0 from FPR
$\approx0.6$ on: a curve hugging the top-left corner (an excellent classifier, AUC close to 1).}

%% ---------------------------------------------------------------- S3.40
\slide{3}{40}{[ROC curve examples]}
\ixh{ROC curve}\ixh{AUC}
\fig{\textbf{Left} (hand-drawn ``xkcd'' style): three ROC curves --- green just above the diagonal (weak), orange in
between, blue highest (jumps to TPR $\approx0.45$ almost immediately, $\approx0.7$ at FPR 0.2, $\approx0.85$ at FPR
0.5). \textbf{Right}: A (purple, AUC 0.923) rises steeply to TPR $\approx0.85$ by FPR $\approx0.12$ then flattens
towards 1; B (blue, AUC 0.794) starts with a jump to $\approx0.17$ at FPR 0 (so it is above A for FPR
$\lesssim0.03$), then rises more gradually ($\approx0.6$ at FPR 0.2, $\approx0.8$ at 0.5); C (red, AUC 0.588) rises
slowly and almost linearly ($\approx0.4$ at FPR 0.2, $\approx0.55$ at 0.4, $\approx0.87$ at 0.8).}
